%% kapa: exam=1 round=32 year=2021 period=2 form=A subject=accounting subject_ko=회계학 seq=28 number=68
%% source: Q-Net 감정평가사 기출문제 (cst00309) artlSeq=5208397 file=extracted/1008/5208397/A형/2021년도 제32회 감정평가사 1차 2교시 A형.pdf
%% conversion: kordoc 4.24.51 -> tools/gen1.py (auto); review=proofread-sonnet; flags=none
\begin{question}{68}
(주)감평은 취득원가 ￦2,500(처분당시 장부금액은 ￦1,500, 원가모형 적용)인 기계장치를 20×1년 초 ￦1,600에 처분하였다. (주)감평은 기계장치 장부금액을 제거하지 않고 처분대가를 잡수익으로 처리하고, 20×1년과 20×2년 각각 취득원가의 10\%를 감가상각비로 계상하였다. 이러한 오류는 20×3년 초 발견되었고, 20×2년도의 장부가 마감되었다면, (주)감평의 20×3년 당기순이익에 미치는 영향은? (단, 상기 오류는 오류의 영향이나 오류의 누적효과를 실무적으로 결정할 수 있으며 중요한 오류에 해당한다.)
\begin{choices}
  \item 영향없음
  \item ￦100 증가
  \item ￦250 증가
  \item ￦500 증가
  \item ￦600 증가
\end{choices}
\end{question}
\answer{1}
